% theme: converse_segment_ratios_and_midpoint_theorem
% 叩き台。解答は未生成。レビュー用図を挿入済み。表3大問、裏4大問。
\MajorQuestionLesson{中点連結定理}
\SetRowSpace{3mm}

%平行線の逆はあまり出ない(大概わかる)ので、中点連結定理の証明問題に絞ります。
%長さ求める系も、普通に相似でいいと思う。
%表面を穴埋め、ウラ面を完全記述に。

%問題の案として、有名なものをいくつか載せます。converse_segment_ratio_proofなものも冒頭にあっていいですが、ページ最後の問題は、この程度の難易度に揃えてください。
%四角形の対角線の中点と、向かい合う辺の中点を結ぶと平行四辺形になることを示す
%四角形の各辺の中点を結ぶと平行四辺形になることを示す
%上下平行の台形があり、左右の中点を結ぶ線の長さは、上底と下底の平均になる


\Needspace{90mm}
\TypeQuestion{次の問いに答えなさい。}{converse_segment_ratio_proof}
\QQSideFig{42mm}{fig/lesson_09_proof_fill}{$\Tri{ABC}$ の辺 $\Lseg{AB},\Lseg{AC}$ 上に点 $\mathrm{D},\mathrm{E}$ をとり、$\Lseg{AD}=6\mathrm{cm},\ \Lseg{DB}=9\mathrm{cm},\ \Lseg{AE}=4\mathrm{cm},\ \Lseg{EC}=6\mathrm{cm}$ とする。$\Lseg{DE}\Parallel\Lseg{BC}$ となることを示す次の証明の空欄をうめなさい。\par\smallskip
\ProofHeading
$\Tri{ADE}$ と $\Tri{ABC}$ において、\par
$\Lseg{AB}=$\Blank{ア}$\mathrm{cm}$、$\Lseg{AC}=$\Blank{イ}$\mathrm{cm}$\par\smallskip
$\Lseg{AD}:\Lseg{AB}=\Lseg{AE}:\Lseg{AC}=$\Blank{ウ}\quad\CondTag{1}\par\smallskip
共通な角であるから、$\Ang{DAE}=\Ang{BAC}$\quad\CondTag{2}\par\smallskip
\CondNo{1}、\CondNo{2}より、\Blank{エ}ので、\par
$\Tri{ADE}\Similar\Tri{ABC}$\par\smallskip
したがって、$\Ang{ADE}=$\Blank{オ}\par\smallskip
\Blank{カ}が等しいので、$\Lseg{DE}\Parallel\Lseg{BC}$
}{}
\QQSideFig{42mm}{fig/lesson_09_proof_write}{$\Tri{ABC}$ の辺 $\Lseg{AB},\Lseg{AC}$ をそれぞれ点 $\mathrm{B},\mathrm{C}$ の側に延長した直線上に点 $\mathrm{D},\mathrm{E}$ をとる。$\Lseg{AB}=8\mathrm{cm},\ \Lseg{BD}=4\mathrm{cm},\ \Lseg{AC}=10\mathrm{cm},\ \Lseg{CE}=5\mathrm{cm}$ のとき、$\Lseg{BC}\Parallel\Lseg{DE}$ となることを証明しなさい。}{}


\AnswerColumnBreak
\LessonQuestionBreak


%証明をまぜつつ、多少辺や周の長さなどを聞く問題があってもいいですね。
\Needspace{90mm}
\TypeQuestion{}{midpoint_theorem_quadrilateral}
\FigProblemBlock{48mm}{fig/lesson_09_quadrilateral}{
\Text{凸四角形 $\mathrm{ABCD}$ の辺 $\Lseg{AB},\Lseg{BC},\Lseg{CD},\Lseg{DA}$ の中点をそれぞれ $\mathrm{E},\mathrm{F},\mathrm{G},\mathrm{H}$ とする。対角線の長さは $\Lseg{AC}=12\mathrm{cm},\ \Lseg{BD}=16\mathrm{cm}$ である。}
}{
\QQ{線分 $\Lseg{EF}$ と対角線 $\Lseg{AC}$ の位置関係を答え、線分 $\Lseg{EF}$ の長さを求めなさい。}{}
\QQ{四角形 $\mathrm{EFGH}$ が平行四辺形であることを証明しなさい。}{}
\QQ{四角形 $\mathrm{EFGH}$ の周の長さを求めなさい。}{}
}


%一直線上に並ぶことを示す、などは避けて。高校っぽい？
%平行四辺形の成立条件を使ったり、二等辺三角形であることを示す、程度はアリです(こういうのと中点連結定理がつながるかは別として、無理に使わなくていいですよ)。

\Needspace{90mm}
\TypeQuestion{}{midpoint_theorem_trapezoid}
\FigProblemBlock{48mm}{fig/lesson_09_trapezoid}{
\Text{$\Lseg{AB}\Parallel\Lseg{CD}$ である台形 $\mathrm{ABCD}$ で、$\Lseg{AB}=10\mathrm{cm},\ \Lseg{CD}=22\mathrm{cm}$ とする。辺 $\Lseg{AD},\Lseg{BC}$ の中点をそれぞれ $\mathrm{E},\mathrm{F}$ とし、対角線 $\Lseg{AC}$ の中点を $\mathrm{M}$ とする。}
}{
\QQ{3点 $\mathrm{E},\mathrm{M},\mathrm{F}$ が一直線上にあることを示し、線分 $\Lseg{EM}$ の長さを求めなさい。}{}
\QQ{線分 $\Lseg{MF}$ の長さを求めなさい。}{}
\QQ{線分 $\Lseg{EF}$ の長さを求めなさい。}{}
}


\Needspace{90mm}
\TypeQuestion{}{midpoint_theorem_medians}
\FigProblemBlock{48mm}{fig/lesson_09_medians}{
\Text{$\Tri{ABC}$ の辺 $\Lseg{BC},\Lseg{CA}$ の中点をそれぞれ $\mathrm{D},\mathrm{E}$ とし、線分 $\Lseg{AD},\Lseg{BE}$ の交点を $\mathrm{G}$ とする。点 $\mathrm{E}$ を通り線分 $\Lseg{AD}$ に平行な直線と、辺 $\Lseg{BC}$ との交点を $\mathrm{F}$ とする。}
}{
\QQ{点 $\mathrm{F}$ が線分 $\Lseg{DC}$ の中点であることを証明しなさい。}{}
\QQ{$\Lseg{EF}:\Lseg{AD}$ を求めなさい。}{}
\QQ{$\Tri{BDG}$ と $\Tri{BFE}$ に注目して、$\Lseg{BG}:\Lseg{GE}$ を求めなさい。}{}
\QQ{$\Lseg{AG}:\Lseg{GD}$ を求めなさい。}{}
}
